% =====================================================================
%  CARNET D'ENTRAÎNEMENT — CHIMIE — FICHE 3 — THÈME 1
%  Particules subatomiques — NIVEAU FACILE — CORRIGÉ
% =====================================================================
\documentclass[11pt,a4paper]{article}

\usepackage[utf8]{inputenc}
\usepackage[T1]{fontenc}
\usepackage{lmodern}
\usepackage[french]{babel}

\usepackage{amsmath,amssymb,mathtools}
\usepackage{xcolor}
\usepackage{colortbl}
\usepackage{tikz}
\usetikzlibrary{arrows.meta,positioning,calc}
\usepackage[most]{tcolorbox}
\usepackage{enumitem}
\usepackage{booktabs}
\usepackage{array}
\usepackage[a4paper,margin=2.1cm,headheight=26pt,headsep=9pt]{geometry}
\usepackage{fancyhdr}
\usepackage[hidelinks]{hyperref}

\definecolor{primary}{RGB}{124,78,140}
\definecolor{primarydark}{RGB}{74,44,92}
\definecolor{defcol}{RGB}{0,114,178}
\definecolor{thmcol}{RGB}{0,140,120}
\definecolor{excol}{RGB}{0,135,90}
\definecolor{attncol}{RGB}{200,40,55}

\newcommand{\nuc}[3]{\prescript{#1}{#2}{\mathrm{#3}}}
\newcommand{\pp}{p^{+}}
\newcommand{\ee}{e^{-}}
\newcommand{\nz}{n^{0}}
\newcommand{\rep}{\textbf{\color{excol}Réponse : }}

\pagestyle{fancy}
\fancyhf{}
\fancyhead[L]{\small\textcolor{primarydark}{\textbf{Fiche 3 — Thème 1}} : Particules subatomiques}
\fancyhead[R]{\small\textcolor{primarydark}{\textsc{Facile — Corrigé}}}
\fancyfoot[C]{\small\thepage}
\renewcommand{\headrulewidth}{0.8pt}
\renewcommand{\headrule}{\hbox to\headwidth{\color{primary}\leaders\hrule height \headrulewidth\hfill}}

\tcbset{
  boxbase/.style={enhanced,breakable,boxrule=0.8pt,arc=2pt,
     left=2.6mm,right=2.6mm,top=1.8mm,bottom=1.8mm,
     fonttitle=\bfseries,coltitle=white,toptitle=0.4mm,bottomtitle=0.4mm},
}
\newtcolorbox{bilan}[1][]{boxbase,colback=thmcol!7,colframe=thmcol,
  title={\textsc{À retenir}\ifx\relax#1\relax\relax\else\textmd{ : \itshape #1}\fi}}

\newcounter{exo}
\newcommand{\cor}[1][]{\stepcounter{exo}\par\medskip%
  \noindent{\color{primary}\rule[-2pt]{3.5pt}{15pt}}\ \textbf{\color{primarydark}\large Exercice \theexo}%
  \ \textmd{\normalsize\color{black!55}— corrigé}%
  \ifx\relax#1\relax\relax\else\ \textmd{\normalsize\color{black!65}\itshape (#1)}\fi\par\nopagebreak\smallskip}

\setlist{leftmargin=1.7em,topsep=2pt,itemsep=3pt}

\begin{document}

\begin{tikzpicture}
\node[rectangle,rounded corners=6pt,fill=primarydark,text=white,
      minimum width=\textwidth,minimum height=1.35cm,align=center,inner sep=6pt]
  {\Large\bfseries Thème 1 — Particules subatomiques\\[2pt]
   \normalsize\mdseries Niveau Facile — Corrigé détaillé};
\end{tikzpicture}

% ---------------------------------------------------------------
\cor[carte d'identité des trois particules]
\begin{center}
\renewcommand{\arraystretch}{1.25}
\begin{tabular}{l c c}
\toprule
\rowcolor{primary!14}
\textbf{Particule} & \textbf{Localisation} & \textbf{Charge} \\
\midrule
Proton $\ \pp$   & noyau              & positive ($+e$) \\
Neutron $\ \nz$  & noyau              & nulle (neutre) \\
Électron $\ \ee$ & nuage électronique & négative ($-e$) \\
\bottomrule
\end{tabular}
\end{center}
Les protons et neutrons (les \textbf{nucléons}) sont dans le noyau ; les électrons gravitent autour, dans
le nuage électronique. \emph{Aucun} électron ne se trouve jamais dans le noyau.

% ---------------------------------------------------------------
\cor[du plus lourd au plus léger]
Ordre : \ $\pp \approx \nz \;>\; \ee$. \\
\rep proton et neutron (masse $\approx 1{,}7\times10^{-27}$ kg chacun) sont les plus lourds et ont une
masse \textbf{quasiment identique} ; l'électron ($\approx 9{,}1\times10^{-31}$ kg) est de très loin le
plus léger (environ $2000$ fois plus léger qu'un nucléon).

% ---------------------------------------------------------------
\cor[vrai ou faux]
\begin{enumerate}[label=\alph*)]
  \item \textbf{\color{attncol}Faux}. Le neutron est \emph{neutre} (charge nulle) — c'est de là que vient
        son nom.
  \item \textbf{\color{attncol}Faux}. L'électron est très léger ($9{,}1\times10^{-31}$ kg), mais sa masse
        n'est \emph{pas} nulle.
  \item \textbf{\color{excol}Vrai}. Le proton porte $+e$, l'électron $-e$ : signes opposés, même valeur
        absolue $e = 1{,}6\times10^{-19}$ C.
  \item \textbf{\color{excol}Vrai}. $m(\ee)/m(\pp)\approx 9{,}1\times10^{-31}/1{,}7\times10^{-27}\approx
        0{,}0005 \approx 1/2000$.
\end{enumerate}

% ---------------------------------------------------------------
\cor[charges élémentaires]
\begin{enumerate}[label=\alph*)]
  \item un proton : $\ +1{,}6\times10^{-19}$ C ;
  \item un électron : $\ -1{,}6\times10^{-19}$ C (charge opposée) ;
  \item cinq protons : $\ 5\times(+1{,}6\times10^{-19}) = +8{,}0\times10^{-19}$ C.
\end{enumerate}

% ---------------------------------------------------------------
\cor[le constituant principal d'un atome]
\rep le \textbf{vide}. \\
L'atome ($\sim10^{-10}$ m) est environ $10^{5}$ fois plus grand que son noyau ($\sim10^{-15}$ m). Entre
le noyau minuscule et les quelques électrons lointains, il n'y a essentiellement rien : l'atome est
\textbf{lacunaire}, son constituant principal en volume est le vide.

\begin{bilan}
\textbf{Charge} : proton $+e$, électron $-e$, neutron $0$, avec $e=1{,}6\times10^{-19}$ C. \quad
\textbf{Masse} : $\pp\approx\nz\gg\ee$ (électron $\sim2000\times$ plus léger). \quad
\textbf{Volume} : l'atome est surtout du \textbf{vide}.
\end{bilan}

\end{document}
